\cotefile{157-4}
\scanpage{196}
\mspage{70}
\deleted{\unsure{Cela} \unsure{résu}} On a bien \unsure{sur|sûr} $\Longleftarrow$, donc il reste à prouver $\Longrightarrow$. Si l'ens \unsure{d'indices} $A$ est fini, cela résulte du lemme 5.5, compte tenu que la finalité \unsure{des} $M_\alpha$ implique leur $\unsure{\pi|R|K}$-cofinalité. On doit donc simplement améliorer 5.5, par le

\uline{Lemme 5.15} \unsure{Soient} \inserted{$\pi$ un cardinal,} $I$ un ens. ordonné filtrant $\pi$-adapté, et soit $(I_\alpha)_{\alpha\in A}$ une famille de parties $I_\alpha$ $\pi$-cofinales (cf 5.8.), telle que $\operatorname{card} A\leq\pi$. Alors $\bigcap I_\alpha$ est $\pi$-cofinale.

Il suffit de prouver qu'elle est cofinale, car la propriété de stabilité par bornes supérieures de parties filtrantes de cardinal $\leq\pi$ est évidente. Quitte à remplacer $I$ par $I/i_0$ et les $I_\alpha$ par leurs traces sur $I/i_0$, la question est simplement de prouver que $\bigcap_\alpha I_\alpha\ne\varnothing$. On va construire un élément $s$ de $\bigcap I_\alpha$ par induction transfinie. Prenons un bon ordre sur $A$, et \deleted{\unsure{construisons} \unsure{une suite|un syst.}} \inserted{\unsure{prouvons} par réc. transfinie \deleted{\unsure{transfinie}}} \deleted{\uline{\unsure{croissante|croissant}} d'éléments $\unsure{i_\alpha|i'_\alpha}$ de $I$, \unsure{telle} que}
% archive sheet number: 195
